\documentclass{article}
\usepackage[utf8]{inputenc}
\usepackage{booktabs}
\usepackage{longtable}
\title{AAP live four-paper fallback}
\begin{document}
\maketitle
\begin{abstract}
AAP live four-paper fallback synthesizes 4 papers across 19 evidence rows using factor structure as the main organizing lens. The synthesis is concentrated in behavioral factors, traffic factors, with recurring methods such as Survey-based Structural Equation Modeling (SEM), Controlled field experiment and task coverage centered on Predicting pedestrian violating crossing behavior intentions, Pedestrian crossing behavior under mobile phone distraction. Across 4 papers, the evidence is most coherently organized around factor structure rather than any single method, and the strongest clusters are behavioral factors, traffic factors, design factors. Future work should prioritize standardized reporting and more explicit cross-study comparison.
\end{abstract}
\noindent\textbf{Keywords:} controlled field experiment, predicting pedestrian violating crossing behavior intentions, pedestrian crossing behavior under mobile phone distraction, identify factors affecting pedestrian red-light running behavior at signalized intersections, 28 college students at signalized intersection in hefei, china, survey-based structural equation modeling (sem), statistical modeling (binomial logistic regression), survey data from 260 respondents in dalian, china
\section{Introduction and Review Scope}
Introduction and Review Scope provides the analytic entry point for this subsection, bringing the discussion into focus around Behavioral Factors, Traffic Factors. Across 3 evidence units, the relevant literature assembles method families such as Statistical Modeling (Binomial Logistic Regression), Discrete Choice Model With Regression Analysis, Survey Based Structural Equation Modeling (Sem), tasks such as Identify Factors Affecting Pedestrian Red Light Running Behavior At Signalized Intersections, Pedestrian Speed Profile Modeling, Crossing Time Estimation, Vehicle Pedestrian Conflict Simulation, Predicting Pedestrian Violating Crossing Behavior Intentions, datasets such as 631 Valid Survey Samples Collected From Six Signalized Intersections In Hefei, China During May June 2015, Video Survey Data From 5 Crosswalks At 3 Intersections In Nagoya City, Japan (Kanayama, Ueda, Fushimi), which defines the comparison set for the discussion. Read this way, the subsection begins from a shared problem space rather than a sequence of isolated paper summaries. \cite{12_2016_zhang_aap,17_2017_chen_aap}

Across the reviewed studies, a recurring pattern emerges in introduction and review scope. Pedestrians who indicate they will run a red light when in hurry are more likely to actually do so (odds ratio = 0.299) The proposed model represents pedestrian travel time distribution more accurately than the constant speed model Collectively, these observations support a subsection-level claim instead of leaving the evidence as a descriptive inventory. \cite{12_2016_zhang_aap,17_2017_chen_aap,08_2016_zhou_aap}

Taken together, the literature on introduction and review scope points to a field-level pattern rather than a single-study result. The evidence clusters most clearly around Behavioral Factors, Traffic Factors. The synthesis therefore clarifies both the main takeaway and the boundary conditions that qualify it. \cite{12_2016_zhang_aap,17_2017_chen_aap}

\section{Factor Taxonomy and Contextual Drivers}
Factor Taxonomy and Contextual Drivers provides the analytic entry point for this subsection, bringing the discussion into focus around Behavioral Factors, Traffic Factors. Across 2 evidence units, the relevant literature assembles method families such as Controlled Field Experiment, Statistical Modeling (Binomial Logistic Regression), tasks such as Pedestrian Crossing Behavior Under Mobile Phone Distraction, Identify Factors Affecting Pedestrian Red Light Running Behavior At Signalized Intersections, datasets such as 28 College Students At Signalized Intersection In Hefei, China, 631 Valid Survey Samples Collected From Six Signalized Intersections In Hefei, China During May June 2015, which defines the comparison set for the discussion. Read this way, the subsection begins from a shared problem space rather than a sequence of isolated paper summaries. \cite{09_2018_jiang_aap,12_2016_zhang_aap}

Across the reviewed studies, a recurring pattern emerges in factor taxonomy and contextual drivers. Text distraction causes the greatest impairment to pedestrian crossing performance, followed by phone conversation distraction and music distraction Pedestrians distracted by texting look left and right less often and switch, distribute and maintain less visual attention on the traffic environment Collectively, these observations support a subsection-level claim instead of leaving the evidence as a descriptive inventory. \cite{09_2018_jiang_aap,12_2016_zhang_aap}

Comparison is most informative where work on factor taxonomy and contextual drivers converges in broad outline but diverges in operational detail. behavioral factors appears in 16 evidence unit(s), with 0 gap signal(s) and 0 contradiction signal(s). traffic factors appears in 14 evidence unit(s), with 0 gap signal(s) and 0 contradiction signal(s). The resulting contrast shows where convergence is substantive and where apparent agreement rests on non-equivalent designs. \cite{09_2018_jiang_aap,12_2016_zhang_aap}

Taken together, the literature on factor taxonomy and contextual drivers points to a field-level pattern rather than a single-study result. The evidence clusters most clearly around Behavioral Factors, Traffic Factors. The synthesis therefore clarifies both the main takeaway and the boundary conditions that qualify it. \cite{09_2018_jiang_aap,12_2016_zhang_aap}

\section{Factor Focus: Behavioral Factors}
Behavioral Factors provides the analytic entry point for this subsection, bringing the discussion into focus around Behavioral Factors, Traffic Factors. Across 2 evidence units, the relevant literature assembles method families such as Survey Based Structural Equation Modeling (Sem), Controlled Field Experiment, tasks such as Predicting Pedestrian Violating Crossing Behavior Intentions, Pedestrian Crossing Behavior Under Mobile Phone Distraction, datasets such as Survey Data From 260 Respondents In Dalian, China, 28 College Students At Signalized Intersection In Hefei, China, which defines the comparison set for the discussion. Read this way, the subsection begins from a shared problem space rather than a sequence of isolated paper summaries. \cite{08_2016_zhou_aap,09_2018_jiang_aap}

Across the reviewed studies, a recurring pattern emerges in behavioral factors. Instrumental attitude was significant in all four models, indicating pedestrians are willing to take risks to cross the road if this behavior is associated with meaningful benefits like saving time or providing convenience Model 4 (containing instrumental attitude, descriptive norm, and conformity tendency) provided the best predictive utility with the best goodness of fit (CFI=0.996, CMIN/DF=1.235, RMSEA=0.030) Collectively, these observations support a subsection-level claim instead of leaving the evidence as a descriptive inventory. \cite{08_2016_zhou_aap,09_2018_jiang_aap}

Taken together, the literature on behavioral factors points to a field-level pattern rather than a single-study result. The evidence clusters most clearly around Behavioral Factors, Traffic Factors. The synthesis therefore clarifies both the main takeaway and the boundary conditions that qualify it. \cite{08_2016_zhou_aap,09_2018_jiang_aap}

\section{Factor Focus: Traffic Factors}
Traffic Factors provides the analytic entry point for this subsection, bringing the discussion into focus around Behavioral Factors, Traffic Factors. Across 2 evidence units, the relevant literature assembles method families such as Survey Based Structural Equation Modeling (Sem), Controlled Field Experiment, tasks such as Predicting Pedestrian Violating Crossing Behavior Intentions, Pedestrian Crossing Behavior Under Mobile Phone Distraction, datasets such as Survey Data From 260 Respondents In Dalian, China, 28 College Students At Signalized Intersection In Hefei, China, which defines the comparison set for the discussion. Read this way, the subsection begins from a shared problem space rather than a sequence of isolated paper summaries. \cite{08_2016_zhou_aap,09_2018_jiang_aap}

Across the reviewed studies, a recurring pattern emerges in traffic factors. Instrumental attitude was significant in all four models, indicating pedestrians are willing to take risks to cross the road if this behavior is associated with meaningful benefits like saving time or providing convenience Model 4 (containing instrumental attitude, descriptive norm, and conformity tendency) provided the best predictive utility with the best goodness of fit (CFI=0.996, CMIN/DF=1.235, RMSEA=0.030) Collectively, these observations support a subsection-level claim instead of leaving the evidence as a descriptive inventory. \cite{08_2016_zhou_aap,09_2018_jiang_aap}

Taken together, the literature on traffic factors points to a field-level pattern rather than a single-study result. The evidence clusters most clearly around Behavioral Factors, Traffic Factors. The synthesis therefore clarifies both the main takeaway and the boundary conditions that qualify it. \cite{08_2016_zhou_aap,09_2018_jiang_aap}

\section{Comparative Factor Evidence and Interactions}
Comparative Factor Evidence and Interactions provides the analytic entry point for this subsection, bringing the discussion into focus around Behavioral Factors, Traffic Factors. Across 2 evidence units, the relevant literature assembles method families such as Survey Based Structural Equation Modeling (Sem), Controlled Field Experiment, tasks such as Predicting Pedestrian Violating Crossing Behavior Intentions, Pedestrian Crossing Behavior Under Mobile Phone Distraction, datasets such as Survey Data From 260 Respondents In Dalian, China, 28 College Students At Signalized Intersection In Hefei, China, which defines the comparison set for the discussion. Read this way, the subsection begins from a shared problem space rather than a sequence of isolated paper summaries. \cite{08_2016_zhou_aap,09_2018_jiang_aap}

Comparison is most informative where work on comparative factor evidence and interactions converges in broad outline but diverges in operational detail. behavioral factors appears in 16 evidence unit(s), with 0 gap signal(s) and 0 contradiction signal(s). traffic factors appears in 14 evidence unit(s), with 0 gap signal(s) and 0 contradiction signal(s). The resulting contrast shows where convergence is substantive and where apparent agreement rests on non-equivalent designs. \cite{08_2016_zhou_aap,09_2018_jiang_aap}

Across the reviewed studies, a recurring pattern emerges in comparative factor evidence and interactions. Instrumental attitude was significant in all four models, indicating pedestrians are willing to take risks to cross the road if this behavior is associated with meaningful benefits like saving time or providing convenience Model 4 (containing instrumental attitude, descriptive norm, and conformity tendency) provided the best predictive utility with the best goodness of fit (CFI=0.996, CMIN/DF=1.235, RMSEA=0.030) Collectively, these observations support a subsection-level claim instead of leaving the evidence as a descriptive inventory. \cite{08_2016_zhou_aap,09_2018_jiang_aap}

Taken together, the literature on comparative factor evidence and interactions points to a field-level pattern rather than a single-study result. The evidence clusters most clearly around Behavioral Factors, Traffic Factors. The synthesis therefore clarifies both the main takeaway and the boundary conditions that qualify it. \cite{08_2016_zhou_aap,09_2018_jiang_aap}

\section{Conclusion}
The field-level picture across 4 papers is that factor structure provide the clearest organizing lens, with the strongest signals concentrated in behavioral factors, traffic factors, design factors. \cite{08_2016_zhou_aap,09_2018_jiang_aap}

Stable conclusions: Across 4 papers, the evidence is most coherently organized around factor structure rather than any single method, and the strongest clusters are behavioral factors, traffic factors, design factors. Methodological diversity is real, but the review still shows recurring families such as Survey-based Structural Equation Modeling (SEM), Controlled field experiment, Statistical modeling (Binomial logistic regression), which suggests the field is comparing similar problems with different operational choices. Reporting remains uneven across tasks and metrics such as Predicting pedestrian violating crossing behavior intentions, Pedestrian crossing behavior under mobile phone distraction and CMIN/df ratio, RMSEA (Root Mean Square Error of Approximation), which limits direct comparison across studies. \cite{08_2016_zhou_aap,09_2018_jiang_aap}

Unresolved tensions: A second tension is that cross-study comparison still depends on inconsistent reporting conventions. \cite{08_2016_zhou_aap,09_2018_jiang_aap}

Research priorities: Future work should prioritize standardized reporting and more explicit cross-study comparison. A third priority is to align methods, metrics, and context descriptors so the literature can be synthesized more defensibly. \cite{08_2016_zhou_aap,09_2018_jiang_aap}

\appendix
\section{Appendix}
\subsection{Appendix Summary}
\noindent The appendix records 4 papers and 19 matrix rows used to generate the review synthesis. The most visible methodological families are Survey-based Structural Equation Modeling (SEM), Controlled field experiment, Statistical modeling (Binomial logistic regression), while recurring tasks include Predicting pedestrian violating crossing behavior intentions, Pedestrian crossing behavior under mobile phone distraction, Identify factors affecting pedestrian red-light running behavior at signalized intersections.

\begin{itemize}
\item \textbf{Papers}: 4
\item \textbf{Matrix rows}: 19
\item \textbf{Verified gaps}: 
\item \textbf{Dominant axis}: factor
\item \textbf{Top methods}: Survey-based Structural Equation Modeling (SEM), Controlled field experiment, Statistical modeling (Binomial logistic regression)
\item \textbf{Top tasks}: Predicting pedestrian violating crossing behavior intentions, Pedestrian crossing behavior under mobile phone distraction, Identify factors affecting pedestrian red-light running behavior at signalized intersections
\item \textbf{Top datasets}: Survey data from 260 respondents in Dalian, China, 28 college students at signalized intersection in Hefei, China, 631 valid survey samples collected from six signalized intersections in Hefei, China during May-June 2015
\end{itemize}

\subsection{Evidence Table}
\begin{longtable}{p{0.14\textwidth}p{0.21\textwidth}p{0.10\textwidth}p{0.17\textwidth}p{0.17\textwidth}p{0.17\textwidth}}
\toprule
Paper & Title & Year & Methods & Tasks & Gap matches \\
\midrule
\endfirsthead
\toprule
Paper & Title & Year & Methods & Tasks & Gap matches \\
\midrule
\endhead
09\_2018\_jiang\_aap & Effects of mobile phone distraction on pedestrians' crossing behavior and visual attention allocation at a signalized intersection: An outdoor experimental study & 2018 & Controlled field experiment & Pedestrian crossing behavior under mobile phone distraction & -- \\
17\_2017\_chen\_aap & Modeling pedestrian crossing speed profiles considering speed change behavior for the safety assessment of signalized intersections & 2017 & Discrete choice model with regression analysis & Pedestrian speed profile modeling, crossing time estimation, vehicle-pedestrian conflict simulation & -- \\
08\_2016\_zhou\_aap & An extension of the theory of planned behavior to predict pedestrians' violating crossing behavior using structural equation modeling & 2016 & Survey-based Structural Equation Modeling (SEM) & Predicting pedestrian violating crossing behavior intentions & -- \\
12\_2016\_zhang\_aap & Exploring factors affecting pedestrians' red-light running behaviors at intersections in China & 2016 & Statistical modeling (Binomial logistic regression) & Identify factors affecting pedestrian red-light running behavior at signalized intersections & -- \\
\bottomrule
\end{longtable}
\end{document}
% BIB START
@article{08_2016_zhou_aap,
  author={Hongmei Zhou; Stephanie Ballon Romero; Xiao Qin},
  title={An extension of the theory of planned behavior to predict pedestrians' violating crossing behavior using structural equation modeling},
  year={2016},
  journal={Accident Analysis and Prevention}
}

@article{09_2018_jiang_aap,
  author={Kang Jiang; Feiyang Ling; Zhongxiang Feng; Changxi Ma; Wesley Kumfer; Chen Shao; Kun Wang},
  title={Effects of mobile phone distraction on pedestrians' crossing behavior and visual attention allocation at a signalized intersection: An outdoor experimental study},
  year={2018},
  journal={Accident Analysis and Prevention}
}

@article{12_2016_zhang_aap,
  author={Weihua Zhang; Kun Wang; Lei Wang; Zhongxiang Feng; Yingjie Du},
  title={Exploring factors affecting pedestrians' red-light running behaviors at intersections in China},
  year={2016},
  journal={Accident Analysis and Prevention}
}

@article{17_2017_chen_aap,
  author={Miho Iryo-Asano; Wael K.M. Alhajyaseen},
  title={Modeling pedestrian crossing speed profiles considering speed change behavior for the safety assessment of signalized intersections},
  year={2017},
  journal={Accident Analysis and Prevention}
}